\documentclass[11pt,a4paper]{article}

\usepackage[utf8]{inputenc}
\usepackage[T1]{fontenc}
\usepackage{amsmath}
\usepackage{amssymb}
\usepackage{graphicx}
\usepackage{booktabs}
\usepackage{tabularx}
\usepackage{xcolor}
\usepackage{hyperref}
\usepackage{listings}
\usepackage{enumitem}
\usepackage{geometry}
\geometry{margin=2.5cm}

% Define colors for code/highlighting
\definecolor{codegreen}{rgb}{0,0.6,0}
\definecolor{codegray}{rgb}{0.5,0.5,0.5}
\definecolor{codepurple}{rgb}{0.58,0,0.82}

% Configuration for listings
\lstset{
    backgroundcolor=\color{white},
    commentstyle=\color{codegreen},
    keywordstyle=\color{magenta},
    numberstyle=\tiny\color{codegray},
    stringstyle=\color{codepurple},
    basicstyle=\ttfamily\footnotesize,
    breakatwhitespace=false,
    breaklines=true,
    captionpos=b,
    keepspaces=true,
    numbers=left,
    numbersep=5pt,
    showspaces=false,
    showstringspaces=false,
    showtabs=false,
    tabsize=2
}

\title{Simulation Scenarios and Parameters}
\author{Muhammad Adeel Zahid}
\date{\today}

\begin{document}

\maketitle

\section{Simulation Scenarios}

The standard simulation scenario is as shown in \ref{fig:system_architecture}. The scenarios to test out are as follows

\begin{itemize}
    \item Multiple Users video streaming directly connected to LEO satellite. Users are static and the satellite is orbiting. The placement of satellite is selected such that in no case the satellite goes outside the field of view of IAB nodes and users
    \item Multiple users, all are connected to a single IAB node
    \item Multiple users, divided equally among 2 IAB nodes and then go on to 4 IAB nodes
    \item User’s coordinates are selected such that they are equal number of users connected to each IAB node. The coordinate placement is discussed in the section below
    \item All users' network and physical layer parameters are kept the same.
    \end{itemize}

The simulation parameters are as follows 

\begin{figure*}[!t]
    \centering
    \includegraphics[width=\textwidth]{Assets/SimulationModel.png}
    \caption{System architecture}
    \label{fig:system_architecture}
\end{figure*}





\section{Video Quality of Experience (QoE) Analysis}

\subsection{Key Metrics}

\begin{itemize}
    \item \textbf{Average Bitrate}: Mean video quality 
    \item \textbf{Minimum Bitrate}: Lowest quality selected
    \item \textbf{Quality Changes}: Number of bitrate switches
    \item \textbf{Total Playback Time}: Effective video playback duration 
    \item \textbf{Interruption Time}: Total rebuffering time
    \item \textbf{Number of Interruptions}: Count of rebuffering events
    \end{itemize}
\subsection{Analysis Approaches}

\subsubsection{QoE Comparison}
Compare QUIC and TCP performance, evaluating:
\begin{itemize}
    \item Average video quality achieved
    \item Quality stability (frequency of switches) to comment on user experience
    \item Cumulative rebuffering time/interrupted time
\end{itemize}

\subsubsection{Bitrate Adaptation/Quality stability Analysis}
To analyze how quickly and accurately the DASH algorithm adapts:
\begin{itemize}
    \item Convergence time to optimal bitrate
    \item Accuracy of throughput estimation (internal mpeg-dash estimation vs end of simulation values
\end{itemize}

\subsection{Visualizations for QoE}
\begin{itemize}
    \item Bitrate over time
    \item Variation of bitrate over time in terms of CDF plot
    \item Quality switch frequency in terms of histogram
    \item Buffer level evolution over time
    \item Playback duration over time
    \item Variation of playback duration interms of CDF plot and box plot
    \item Interruption duration over time
    \item Variation of interruption duration over time
\end{itemize}

\section{Network Performance Analysis}

\subsection{Throughput Analysis}

\subsubsection{Key Metrics}
\begin{itemize}
    \item Throughput over time
    \item Variation of throughput over time in terms of a CDF plot

\end{itemize}

\subsubsection{Fairness Analysis}

Evaluate fairness using Jain's Fairness Index:
\begin{equation}
    J = \frac{\left(\sum_{i=1}^{n} x_i\right)^2}{n \times \sum_{i=1}^{n} x_i^2}
\end{equation}

where $x_i$ is the throughput of user $i$ and $n$ is the number of users. The index ranges from $\frac{1}{n}$ (worst case) to 1 (perfect fairness).

\textbf{Analysis Dimensions:}
\begin{itemize}
    \item Compare fairness across QUIC and TCP
    \item Analyze fairness with different numbers of UEs: 2, 5, 10, 20
    \item Evaluate fairness under different IAB topologies (1 IAB, 2 IAB ...)
\end{itemize}

\subsubsection{Throughput Efficiency}
\begin{itemize}
    \item \textbf{Goodput}: Application-level throughput (including retransmissions)
        \begin{equation}
            \text{Goodput} = \frac{\text{Successfully delivered data}}{\text{Total time}}
        \end{equation}
    \item \textbf{Efficiency Ratio}: $\text{Efficiency} = \frac{\text{Goodput}}{\text{Throughput}}$
    \item \textbf{Protocol overhead comparison} (QUIC vs TCP headers)
\end{itemize}

\subsubsection{Throughput Stability}
Measure throughput variability:
\begin{itemize}
    \item Coefficient of variation: $CV = \frac{\sigma}{\mu}$
    \item Throughput variance over time
\end{itemize}

\subsection{Latency Analysis}

\subsubsection{Key Metrics}
\begin{itemize}
    \item Round-Trip Time (RTT) from QUIC traces
    \item Segment fetch time (from DASH client)
    \item End-to-end latency
\end{itemize}

\subsubsection{Latency Distribution Analysis}
\begin{itemize}
    \item Statistical plot interms of CDF and box plot
    \item Latency vs distance from satellite/IAB
    \item Jitter analysis: variability in latency
        \begin{equation}
            \text{Jitter} = \sqrt{\frac{1}{n-1}\sum_{i=1}^{n}(L_i - \bar{L})^2}
        \end{equation}
        where $L_i$ is the latency of packet $i$ and $\bar{L}$ is the mean latency
\end{itemize}

\subsubsection{Impact on QoE}
\begin{itemize}
    \item Correlation between latency and rebuffering events
    \item Segment fetch time impact on adaptation decisions
\end{itemize}

\subsubsection{Protocol Comparison}
\begin{itemize}
    \item QUIC connection establishment time vs TCP three-way handshake
    \item RTT smoothing behavior comparison
\end{itemize}

\subsection{Congestion Control Behavior}
\textbf{Key Metrics:}
\begin{itemize}
    \item Congestion window evolution over time
    \item Window growth rate during slow start
    \item Window reduction on loss events
\end{itemize}

\subsubsection{Recovery Mechanisms}
\begin{itemize}
    \item Effectiveness of packet recovery from out-of-order packets (To-be-studied in detail since QUIC does not recover at all with out-of-order packet loss)
    \item Time to recover from losses
    \item Throughput degradation during recovery
    \item Recovery success rate
\end{itemize}

\section{IAB Topology Impact Evaluation}

Evaluate performance (throughput and SINR) with varying configurations:
\begin{itemize}
    \item Number of IAB nodes: 1, 2, 3, 4
    \item Number of UEs
    \item UEs per IAB ratio
    \item System bottleneck identification (satellite link vs access link)
    \item Mobility of satellite and user
\end{itemize}


\end{document}
